
\section{Covers charged to a single sausage mesh}
\label{sec:dense-cover}

We prove finite expected $h$-measure of the entire trace in each bounded
region. Select disjoint disks carrying enough normalized sausage area,
all measured at one deterministic mesh. Their enlarged disks have a
uniform expected covering cost; the visited cells they miss have a cost
that tends to zero. Keeping that bound independent of the time and
interior cutoffs gives the whole-time spatial conclusion.

Use the parameters from \eqref{eq:parameters} and the gauge
from \eqref{eq:gauge}. Put
\(\Gamma=\gamma([0,\infty))\).  For \(v>0\), define the random measure
\begin{equation}
 \mu^v(A)=v^{-p}\int_A
       {\bf1}_{\{\operatorname{dist}(w,\Gamma)<v\}}\,dA(w).
 \label{eq:cover-sausage}
\end{equation}
The sausage is taken in the whole plane.  We also use \(\mu_t^v\),
defined by replacing \(\Gamma\) with \(\gamma([0,t])\).  In particular,
\(\mu_t^v\le\mu^v\).

\paragraph{The stopped estimates used by the cover.}
For an interior marked point \(z\), let \(R_t(z)\) be its conformal
radius in the surviving domain, and let
\[
 M_t(z)=R_t(z)^{-p}\sin^{a_0}\theta_t(z),\qquad
 M_0(z)=(2\operatorname{Im}z)^{-p}
        \left(\frac{\operatorname{Im}z}{|z|}\right)^{a_0}.
\]
Let \(\mathbb P_z^*\) denote the finite-scale tilted laws.  Lemmas~\ref{ll:finite-tilt}, \ref{ll:inverse-angle}, and~\ref{ll:approach}
give consistency at finite conformal-radius stops, conditional change
of law, and the following two bounds.
If \(\widehat S\) is one unit of the clock \(-\log R_t(z)\)
after a stop \(S\), then
\begin{equation}
 \mathbb E_z^*\!\left[
      \sin^{-a_0}\theta_{\widehat S}(z)\mid\mathcal F_S\right]\le C.
 \label{eq:cover-angle}
\end{equation}
If \(S\) is a stop at which \(z\) is alive and
\(\rho\le e^{-1}R_S(z)\), then
\begin{equation}
 \mathbb P\{\tau_\rho(z)<\infty\mid\mathcal F_S\}
      \le C\rho^p M_S(z),
 \label{eq:cover-one-point}
\end{equation}
where \(\tau_\rho\) here denotes the subsequent first stop with
conformal radius \(\rho\).  Under the tilted law every finite clock
level is reached.  The finite-stop identity takes the form
\begin{equation}
 \mathbb E\!\left[
       {\bf1}_{\{S<\infty\}}M_S(z)X\right]
        =M_0(z)\mathbb E_z^*[X]
 \label{eq:cover-finite-tilt}
\end{equation}
for nonnegative variables \(X\) measurable at a fixed conformal-radius
stop \(S\), with the killed density assigned value zero.

Proposition~\ref{rc:dense} supplies the following precise uniformity.  Fix a
positive angle cutoff.  There are constants \(b,c,C,C_2>0\), \(C_1>1\),
and an integer \(m_0\) such that, for every integer \(m\ge m_0\)
and every radius \(r>0\), there is a deterministic
\(\varepsilon(r,m)=r\delta_m>0\) with this property.  At any state with
\(R(z)=r\) and angle sine above the cutoff, and for every prescribed
\(0<v\le\varepsilon(r,m)\), the tilted conditional probability is
at least \(c\exp(-Cm^2)\) that, before the radius falls to
\(r/C_1\), the future trace generates in \(B(z,C_2r)\) a
\(v\)-sausage of normalized area at least \(b r^d m^p\).
The constants and the cutoff \(\varepsilon(r,m)\) are uniform in
the entire stopped state.  The event is observable by that clock
endpoint.  This formulation permits a single deterministic physical
mesh to be used at all the tests below.

We also use the elementary visit geometry associated with the
surviving-domain convention.  If \(\overline B(z,Lr)\subset\mathbb H\)
initially, with \(L\) a sufficiently large absolute constant, then a
visit to the square of side \(r\) centered at \(z\) requires a finite
live stop at conformal radius \(Lr\).  Indeed, while \(z\) survives,
Koebe gives
\(R_t(z)\le4\operatorname{dist}(z,\partial D_t)\).
If \(z\) is swallowed without the trace first meeting
\(\overline B(z,2r)\), that ball is in the swallowed component and
the future trace cannot enter it.  Otherwise, just before the first
such meeting the live radius has decreased to at most \(8r\);
continuity gives the earlier radius crossing, for example with
\(L=16\).  The same argument shows that a trace visit within
distance \(v\) of \(z\) forces a live radius crossing at \(16v\),
provided \(\overline B(z,16v)\subset\mathbb H\) initially.
These statements concern the trace, rather
than the filled hull, and apply also when \(\kappa>4\).

\begin{lemma}[Uniform bounded-region charge]
\label{lem:cover-charge}
For every \(0<R<\infty\), there is \(C_R<\infty\), depending only on
\(R\) and \(\kappa\), such that
\[
 \sup_{0<v\le1}\mathbb E\mu^v(B(0,R))\le C_R.
\]
\end{lemma}

\begin{proof}
Choose a fixed \(A>8e\).  If \(y=\operatorname{Im}w>Av\),
the visit geometry and \eqref{eq:cover-one-point}, with a fixed
multiplicative slack in the target radius, give
\[
 \mathbb P\{\operatorname{dist}(w,\Gamma)<v\}
       \le C v^p M_0(w).
\]
For example, increasing \(A\) if necessary makes the target
\(16v\) at least one clock unit below the initial radius \(2y\).
Since \(\Gamma\subset\overline{\mathbb H}\), only the layer
\(-v<y\le Av\) contributes outside this range.  Its area in
\(B(0,R)\) is at most \(C_Rv\).  Tonelli's theorem therefore gives
\begin{align*}
 \mathbb E\mu^v(B(0,R))
 &\le C_Rv^{1-p}
       +C\int_{B(0,R)\cap\mathbb H}M_0(w)\,dA(w)\\
 &\le C_Rv^{1-p}
       +C\int_{-R}^{R}\int_0^R(2y)^{-p}\,dy\,dx
 \le C_R.
\end{align*}
Here \(p<1\) is exactly what makes the boundary integral finite.
The estimate concerns the whole trace at all times.
\end{proof}

\begin{proposition}[The upper Hausdorff bound]
\label{prop:dense-cover-upper}
For every \(0<R<\infty\),
\begin{equation}
 \mathbb E\mathcal H^h
       \bigl(\Gamma\cap\overline B(0,R)\bigr)\le C_R<\infty.
 \label{eq:cover-all-time-expectation}
\end{equation}
Consequently, on one probability-one event,
\[
 \mathcal H^h\bigl(\Gamma\cap([-m,m]+i[0,m])\bigr)<\infty
       \qquad(m=1,2,\ldots).
\]
On that event every compact capacity-time segment has finite
\(\mathcal H^h\)-measure.
\end{proposition}

\begin{proof}
First fix a deterministic compact set
\(K\subset\mathbb H\cap\overline B(0,R)\).
Choose \(T_0>1+\log C_1\), and put
\[
 r_j=e^{-T_0j},\qquad
 J_0=\lfloor n/3\rfloor,\qquad J_1=\lfloor2n/3\rfloor,
 \qquad m_n=\lfloor c_0\sqrt{\log n}\rfloor.
\]
The constant \(c_0>0\) will be chosen below.  We omit finitely many
initial \(n\), so that \(m_n\ge m_0\) and every radius tested is
smaller than the initial conformal radius at the grid centers used.

\emph{One deterministic mesh.}
Choose a deterministic sequence \(v_n>0\) such that \(v_n\to0\)
and
\begin{equation}
 v_n\le\min_{J_0\le j\le J_1}
              \varepsilon(e^{-1}r_j,m_n).
 \label{eq:cover-mesh-choice}
\end{equation}
For example, take half this finite minimum and also require
\(v_n\le2^{-n}\) and \(v_n\le v_{n-1}/2\).  This choice is
deterministic because the dense-visit cutoff is uniform in the
state.  Write \(\mu^{(n)}=\mu^{v_n}\).

Use the usual grid of closed squares of side \(r_n\), and keep
the finitely many squares meeting \(K\).  A center \(z\) of such
a square has initial height bounded below once \(n\) is large.
For the argument at this one center, use \(\mathbb P_z^*\) only
through the finite clock levels required.  Let \(S_j(z)\) be the
first live stop with \(R(z)=r_j\), and let \(S'_j(z)\) be the
stop one clock unit later, so that \(R(z)=e^{-1}r_j\).
All these stops are finite under \(\mathbb P_z^*\).

By \eqref{eq:cover-angle}, a fixed angle cutoff can be chosen
so that the conditional probability of a good angle at \(S'_j\),
given the complete history through \(S_j\), is at least \(1/2\).
The dense-visit input and \eqref{eq:cover-mesh-choice} then give
a conditional success probability at least
\[
 c\exp(-Cm_n^2).
\]
On success, the prescribed mass has been generated before
\(S_{j+1}\), since \(T_0>1+\log C_1\).  Choose a fixed
\(C_3\ge1\) so that all these mass tests lie in
\[
 D_j(z)=\overline B(z,C_3r_j).
\]
There is a constant \(b_1>0\), independent of \(n,j,z\), such that
success implies
\begin{equation}
 \mu_{S_{j+1}(z)}^{v_n}(D_j(z))
     \ge b_1 r_j^d(\log n)^{p/2}.
 \label{eq:cover-adapted-success}
\end{equation}
This uses \(m_n\ge(c_0/2)\sqrt{\log n}\) for large \(n\).
Let \(A_{n,j}(z)\) be the event that \(S_{j+1}(z)<\infty\) and
\eqref{eq:cover-adapted-success} holds. Under the tilted law all these
endpoint stops are finite. The event \(A_{n,j}(z)\) is measurable at
\(S_{j+1}(z)\), because its sausage uses only the trace through that stop.
Choose \(c_0\) so small that \(Cc_0^2<1/2\).  For all sufficiently
large \(n\), the dense-visit lower bound therefore gives
\[
 \mathbb P_z^*\!\left(A_{n,j}(z)\mid\mathcal F_{S_j(z)}\right)
       \ge n^{-1/2}.
\]

Put
\[
 F_n(z)=\bigcap_{j=J_0}^{J_1}A_{n,j}(z)^c.
\]
Every preceding test is observable when the next one starts.
Iterated conditional expectations give
\begin{equation}
 \mathbb P_z^*(F_n(z))
 \le(1-n^{-1/2})^{J_1-J_0+1}
 \le C\exp(-c\sqrt n).
 \label{eq:cover-failure}
\end{equation}
No independence is required between tests or between centers.

\emph{Returning to the original law.}
Call \(D_j(z)\) a candidate if
\begin{equation}
 \mu^{(n)}(D_j(z))\ge b_1r_j^d(\log n)^{p/2}.
 \label{eq:cover-candidate}
\end{equation}
This is a terminal-sausage test. A cell visited by the trace and
having no candidate must satisfy the adapted failure event \(F_n(z)\):
at each finite endpoint, \(\mu_{S_{j+1}}^{v_n}\le\mu^{(n)}\).
By the visit geometry, the visit also requires the
live stop \(\tau_{Lr_n}(z)<\infty\), for the fixed constant
\(L\) above.  For large \(n\), that stop is after
\(S=S_{J_1+1}(z)\) with a gap of at least one clock unit.
Thus on the event on the left below all the test endpoints are finite,
and \(F_n(z)\) is measurable at \(S\).
Applying \eqref{eq:cover-one-point}, followed by
\eqref{eq:cover-finite-tilt}, gives
\begin{align}
 &\mathbb P\{\text{the cell at }z\text{ is visited and has no candidate}\}
       \notag\\
 &\quad\le C r_n^p\,
       \mathbb E\!\left[
         {\bf1}_{\{S<\infty\}}M_S(z){\bf1}_{F_n(z)}\right]
       \notag\\
 &\quad= C r_n^p M_0(z)\mathbb P_z^*(F_n(z))
       \le C r_n^p M_0(z)e^{-c\sqrt n}.
 \label{eq:cover-missed-cell}
\end{align}
The full-sausage candidate information enters only through the
inclusion in the adapted failure event.  All changes of law
occur at finite conformal-radius levels.

\emph{The finite cover and its charge.}
Order the finite collection of candidate disks by decreasing
radius, using a fixed deterministic order to break ties.
Select the first, discard all disks meeting it, and repeat.
The selected closed disks are pairwise disjoint.  A discarded
disk has no greater radius than the selected disk meeting it,
so it is contained in the concentric triple of that selected
disk.  For large \(n\), every cell is contained in each of its
candidate disks, since \(r_n/r_{J_1}\to0\).  Consequently the
triples of the selected disks cover all cells having a candidate.

Uniformly for \(J_0\le j\le J_1\), formula
\eqref{eq:gauge} implies
\[
 h(6C_3r_j)\le C r_j^d(\log n)^{p/2}.
\]
The qualification \eqref{eq:cover-candidate} and disjointness
therefore show that the total cost of these triples is at most
\begin{equation}
 C\sum_{\text{selected }D}\mu^{(n)}(D)
       \le C\mu^{(n)}(B(0,R+1)).
 \label{eq:cover-selected-charge}
\end{equation}
All the candidate disks lie in this fixed larger disk once
\(n\) is large.  No limiting measure or convergence of sausages
is needed in this estimate: every selected disk is charged to
the same measure \(\mu^{(n)}\).

Add every retained fine cell that is visited and has no candidate.
These cells, together with the selected triples, form a finite
cover of \(\Gamma\cap K\).  There are at most \(C_Rr_n^{-2}\)
retained cells; their centers stay a positive distance from the
real line, so \(M_0(z)\le C_K\).  Since
\(h(\sqrt2r_n)\le Cr_n^d(\log n)^{p/2}\),
\eqref{eq:cover-missed-cell} yields an expected added cost at most
\begin{equation}
 C_Kr_n^{-2}r_n^d(\log n)^{p/2}r_n^p e^{-c\sqrt n}
   =C_K(\log n)^{p/2}e^{-c\sqrt n}\longrightarrow0.
 \label{eq:cover-cleanup}
\end{equation}
Here \(d+p=2\).  The maximal diameter of the cover is at most
\[
 \delta_n=\max\{6C_3r_{J_0},\sqrt2r_n\}\longrightarrow0.
\]
If \(Z_n(K)\) is its total cost, Lemma~\ref{lem:cover-charge}
and \eqref{eq:cover-selected-charge}--\eqref{eq:cover-cleanup}
show that
\begin{equation}
 \limsup_{n\to\infty}\mathbb E Z_n(K)\le C_R.
 \label{eq:cover-cost-expectation}
\end{equation}
The constant \(C_R\) does not depend on the distance of \(K\)
from the real line.  Only the vanishing error and the discarded
initial values of \(n\) may depend on \(K\).

\emph{Measurability and the Hausdorff limit.}
For fixed \(v\), the distance to the whole trace is the infimum
over nonnegative rational times of \(|w-\gamma(t)|\).  Thus all
the sausage integrals are measurable, by Tonelli's theorem.
There are finitely many candidate indicators and the greedy
order is deterministic.  A fixed closed cell meets the trace
if and only if it meets \(\gamma([0,N])\) for some integer
\(N\); each event is measurable by continuity of \(\gamma\).
It follows that \(Z_n(K)\) is measurable.

For completeness, \(\mathcal H^h(\gamma([0,T])\cap K)\) is
also measurable for deterministic \(T<\infty\) and compact
\(K\).  In defining the Hausdorff limit of this random compact
set, one may use finite covers whose members are finite unions
of rational open balls.  There are only countably many such
covers.  Their covering events are measurable: for an open
union \(U\), the event
\(\gamma([0,T])\cap K\subset U\) is the event that the compact
trace avoids the deterministic compact set \(K\setminus U\).
Arbitrarily small open enlargements, using continuity of \(h\),
show that these countable covering infima have the same limit
as the usual Hausdorff contents.  This proves measurability.

For every fixed \(T\), the covers above cover
\(\gamma([0,T])\cap K\).  Therefore, pathwise,
\[
 \mathcal H^h(\gamma([0,T])\cap K)
       \le\liminf_{n\to\infty}Z_n(K).
\]
Indeed every sufficiently late cover is admissible at any
fixed positive diameter cutoff.  Fatou's lemma and
\eqref{eq:cover-cost-expectation} give
\begin{equation}
 \mathbb E\mathcal H^h(\gamma([0,T])\cap K)\le C_R,
 \label{eq:cover-compact-expectation}
\end{equation}
uniformly in both \(T\) and \(K\subset\mathbb H\cap\overline B(0,R)\).

\emph{The entire trace in a bounded region.}
To retain this uniformity when approaching both infinite time
and the real boundary, set
\[
 K_q=\overline B(0,R)\cap\{\operatorname{Im}z\ge1/q\},
 \qquad E_q=\gamma([0,q])\cap K_q
       \quad(q=1,2,\ldots).
\]
The sets \(E_q\) are increasing compact sets, and
\[
 \bigcup_{q\ge1}E_q=\Gamma\cap\overline B(0,R)\cap\mathbb H.
\]
Gauge Hausdorff measure is a measure on Borel sets.  Its
continuity from below and then monotone convergence, applied
to \eqref{eq:cover-compact-expectation}, give
\[
 \mathbb E\mathcal H^h
       (\Gamma\cap\overline B(0,R)\cap\mathbb H)
    =\lim_{q\to\infty}\mathbb E\mathcal H^h(E_q)
    \le C_R.
\]
Finally, a bounded real interval is covered by \(O_R(r^{-1})\)
intervals of length \(r\), whose total cost is
\[
 O_R\!\left(r^{d-1}(\log\log(1/r))^{p/2}\right)\longrightarrow0,
\]
because \(d>1\).  The real-axis part therefore has zero
\(\mathcal H^h\)-measure.  This proves
\eqref{eq:cover-all-time-expectation}.

Apply the result at integer radii and intersect the resulting
countably many probability-one events.  Every box in the
statement is contained in one of these disks.  Every compact
time segment has bounded image by continuity of \(\gamma\),
so it too lies in such a disk.  This proves both assertions
simultaneously.
\end{proof}

\begin{proof}[Proof of Theorem~\ref{thm:main}]
Section~\ref{retain:positivity} gives positivity simultaneously on all real
$0<s<t<\infty$ for the fixed parameter $\kappa$.
Proposition~\ref{prop:dense-cover-upper} gives the all-time bounded-box
estimate and finite measure on every compact time segment on another
probability-one event. Their intersection proves all assertions.
\end{proof}
